\documentclass{article}
\usepackage[utf8]{inputenc}
\usepackage[T1]{fontenc}
\usepackage{booktabs}
\usepackage{siunitx}
\usepackage{xintexpr}
\usepackage{pgffor} % Para los bucles \foreach

\sisetup{
	round-mode = places,
	round-precision = 4,
}

% Definición de la función (mantenemos la versión float para n=20)
\newcommand{\trxn}[2]{%
	\xintthefloatexpr
	((#1)/(#2)) * (
	(exp(0) + exp(-(#1)^2))/2 + 
	add(exp(-(i*(#1)/(#2))^2), i=1..#2-1)
	)
	\relax
}

\begin{document}
	
	\section*{Tabla Automatizada con \texttt{pgffor}}
	
	% 1. Definimos una macro para ir guardando las filas
	\newcommand{\cuerpotabla}{}
	
	% 2. Construimos las filas con un bucle anidado
	\foreach \i in {0, 0.5, 1.0, 1.5, 2.0} {
		% Guardamos el valor de x en la primera columna
		\xdef\filatemp{\i} 
		
		% Bucle para las columnas (n = 8, 10, 14, 16, 20)
		\foreach \n in {8, 10, 14, 16, 20} {
			% Calculamos el valor y lo añadimos a la fila con un &
			\xdef\valor{\trxn{\i}{\n}}
			\xdef\filatemp{\filatemp & \valor}
		}
		
		% Añadimos la fila completa al cuerpo de la tabla y cerramos con \\
		\xdef\cuerpotabla{\cuerpotabla \filatemp \\}
	}
	
	\begin{table}[h]
		\centering
		\caption{Resultados automatizados para $n=\{8, 10, 14, 16, 20\}$}
		\vspace{1em}
		\begin{tabular}{c *{5}{S}}
			\toprule
			$x$ & {$n=8$} & {$n=10$} & {$n=14$} & {$n=16$} & {$n=20$} \\
			\midrule
			\cuerpotabla
			\bottomrule
		\end{tabular}
	\end{table}
	
\end{document}